% \pdfoutput=1
\documentclass{article}
\usepackage[table]{xcolor}
\usepackage{iclr2025_conference,times}
\usepackage[utf8]{inputenc} % allow utf-8 input
\usepackage[T1]{fontenc}    % use 8-bit T1 fonts
\usepackage{hyperref}       % hyperlinks
\usepackage{url}            % simple URL typesetting
\usepackage{booktabs}       % professional-quality tables
\usepackage{amsfonts}       % blackboard math symbols
\usepackage{nicefrac}       % compact symbols for 1/2, etc.
\usepackage{microtype}      % microtypography
\usepackage{amsmath}
\usepackage{amssymb}
\usepackage{array}
\usepackage{pifont}
\usepackage{threeparttable}
\usepackage{graphicx} % for \scalebox
% \usepackage{subcaption}
\usepackage{array} % for \tabcolsep
\usepackage{makecell}
\usepackage{algorithm} 
\usepackage{algpseudocode} 
\usepackage{multirow}
\usepackage{listings}
\usepackage{enumitem}
\usepackage{subfig}
% \usepackage{ulem}
\usepackage{float}
\definecolor{customblue}{HTML}{ADD8E6} 

\usepackage{tcolorbox}

\usepackage{footmisc} 
% Define your background and placeholder colors
\definecolor{beigebg}{RGB}{253,249,238} % Light beige
\definecolor{argcolor}{RGB}{70,130,180} % Steel blue for arg-inputs

% Define a tcolorbox for the text with a beige background
\tcbset{
  mypromptstyle/.style={
    colback=beigebg,
    colframe=beigebg,
    boxrule=0pt,
    sharp corners,
    fontupper=\ttfamily\scriptsize,
    before=\vspace{0.5em},
    after=\vspace{0.5em},
  }
}

\newcommand{\cmark}{\textcolor{green}{\ding{51}}} % green check mark
\newcommand{\xmark}{\textcolor{red}{\ding{55}}} % red cross mark

\newcommand{\etal}{\textit{et al}. }
\newcommand{\etc}{\textit{etc}.}
\newcommand{\eg}{\textit{e.g.}}
\newcommand{\aka}{\textit{a.k.a}}
\newcommand{\ie}{\emph{i.e.}}

\usepackage{authblk}
\newcommand{\equalcontrib}{\textsuperscript{\textdagger}}
\newcommand{\correspondauthor}{\textsuperscript{*}}

\title{\Large{ToolACE: Winning the Points of LLM Function Calling}}


 \author[1]{Weiwen Liu\equalcontrib}
 \author[3]{Xu Huang\equalcontrib}
 \author[2]{Xingshan Zeng\equalcontrib}
 \author[2]{Xinlong Hao}
 \author[2]{Shuai Yu}
 \author[2]{Dexun Li}
 \author[2]{Shuai Wang}
 \author[2]{Weinan Gan}
 \author[2]{Zhengying Liu}
 \author[5]{Yuanqing Yu}
 \author[6]{Zezhong Wang}
 \author[4]{Yuxian Wang}
 \author[4]{Wu Ning}
 \author[4]{Yutai Hou}
 \author[2]{Bin Wang}
 \author[2]{Chuhan Wu\correspondauthor}
 \author[2]{Xinzhi Wang}
 \author[2]{Yong Liu}
 \author[2]{Yasheng Wang\correspondauthor}
 \author[4]{Duyu Tang}
 \author[4]{Dandan Tu}
 \author[2]{Lifeng Shang}
 \author[2]{Xin Jiang}
 \author[2]{Ruiming Tang\correspondauthor}
 \author[3]{Defu Lian\correspondauthor}
 \author[2]{Qun Liu}
 \author[3]{Enhong Chen}

 \affil[1]{Shanghai Jiao Tong University}
 \affil[2]{Huawei Noah's Ark Lab}
 \affil[3]{University of Science and Technology of China}
 \affil[4]{Huawei Technologies Co., Ltd}
 \affil[5]{Tsinghua University}
 \affil[6]{The Chinese University of Hong Kong \protect\\ \ wwliu@sjtu.edu.cn,\,zeng.xingshan@huawei.com,\,xuhuangcs@mail.ustc.edu.cn}


\iclrfinalcopy
 
\begin{document}
{
\renewcommand{\thefootnote}{\textdagger} \footnotetext[1]{Equal Contributions.\,\,\textsuperscript{*}Corresponding authors.} 
\renewcommand{\thefootnote}{\fnsymbol{footnote}}
}


\maketitle


\begin{abstract}
  Function calling significantly extends the application boundary of large language models (LLMs), where high-quality and diverse training data is critical for unlocking this capability. However, collecting and annotating real function-calling data is challenging, while synthetic data from existing pipelines often lack coverage and accuracy. In this paper, we present ToolACE, an automatic agentic pipeline designed to generate accurate, complex, and diverse tool-learning data, specifically tailored to the capabilities of LLMs. ToolACE leverages a novel self-evolution synthesis process to curate a comprehensive API pool of 26,507 diverse APIs. Dialogs are further generated through the interplay among multiple agents, under the guidance of a complexity evaluator. To ensure data accuracy, we implement a dual-layer verification system combining rule-based and model-based checks. We demonstrate that models trained on our synthesized data---even with only 8B parameters---achieve state-of-the-art performance, comparable to the latest GPT-4 models. Our model and a subset of the data are publicly available at \url{https://huggingface.co/Team-ACE}.

\end{abstract}

% The number of these API parameters can range from zero to dozens, and the application domains of these APIs span a wide range of businesses and industries~\cite{ye2024tooleyes}.
% For instance, we observe that real-world API parameters are not limited to simple strings or numbers but also include lists, dictionaries, and nested types such as lists of lists or lists of dictionaries.

\section{Introduction}
Equipping Large Language Models (LLMs) with external tools has significantly enhanced the capability of AI Agents to solve complex real-world tasks~\cite{huang2024planning,qin2023toolllm,qu2024tool}. The integration of function calling enables LLMs to access up-to-date information, perform delicate computations, and utilize third-party services, thereby unlocking a wide range of potential applications across various fields, \eg, workflow automation~\cite{zhong2023llm4eda}, financial reporting~\cite{theuma2024equipping}, and travel planning~\cite{hao2024large}. 

Function calls in real-world applications are often diverse and complex, driven by the varied functionalities of APIs\footnote{In this paper, APIs, tools, functions, and plugins are used interchangeably.} and the broad range of tasks~\cite{qin2023toolllm}. APIs often undergo rapid updates to meet diverse user needs, necessitating models capable of robust zero-shot generalization. Additionally, users' requirements can be complex or ambiguous, leading to scenarios where multiple tools are employed in a parallel or dependent manner, or require multi-turn interactions. This highlights the importance of managing intricate instructions and accommodating various function-calling scenarios.


Despite these challenges, current tool-augmented LLMs primarily focus on simple function-calling tasks with limited diversity and complexity~\cite{qu2024tool}. They mainly rely on existing public APIs for task construction, which restricts their zero-shot capabilities and applicability to single-turn queries, neglecting more complex scenarios such as dependent or multi-turn interactions~\cite{qin2023toolllm,tang2023toolalpaca,liu2024apigen}. Table~\ref{table:comparison} provides an overview of the data statistics used in these representative tool-augmented LLMs. Moreover, executions of function calls demand precise API selection and parameter configuration, which are highly dependent on the quality and accuracy of underlying data. 
As data becomes increasingly diverse and complex, generating accurate samples with simple pipelines introduced by the existing work becomes significantly more challenging.


% Driven by practical applications, function calls can be diverse and complex due to the functionalities of real-world APIs~\cite{qin2023toolllm}\footnote{In this paper, APIs, tools, functions, and plugins are interchangeable.}
% Real-world API parameters often extend beyond simple strings or numbers to include complex types like lists, dictionaries, and nested parameters.
% Moreover, a single API is often insufficient to complete a task; instead, multiple tools must be used together to perform real-world tasks~\cite{huang2024planning}.
% The input of one API may even depend on the output of another~\cite{qin2023toolllm}, making function calls even more complicated and challenging.

% Function calls for specific tasks or user queries can generally be grouped into three categories: \textit{single function calls}, \textit{parallel function calls}, and \textit{dependent function calls}. In a single function call, the LLM selects and executes one function to accomplish the user's task. For parallel function calls, the LLM performs multiple independent function calls simultaneously within one turn. Whereas dependent function calls involve the LLM making a series of sequential calls, with each subsequent call relying on the output of the previous ones.

% However, current tool-augmented LLMs primarily focus on simple function-calling scenarios with limited diversity and complexity~\cite{qu2024tool}. 
% Table~\ref{table:comparison} provides an overview of the data statistics utilized in these representative tool-augmented LLMs.
% Although satisfactory performance has been achieved on the single function calling that executes one API in one turn, the capabilities of parallel and dependent function calls are largely overlooked. Furthermore, constrained API domains, simplistic parameter types, and uniform data formats may hinder the applicability of function calling to more complex real-world tasks.
% Additionally, executing real-world function calls requires precise API selection and accurate parameter configuration, both of which are closely tied to the accuracy of the underlying data. However, ensuring data accuracy for tool learning remains a challenging problem, particularly when dealing with diverse and complex tasks.

\begin{table}[t!]
\centering
\small
\begin{threeparttable}
\caption{Comparison of ToolACE with other representative tool-augmented LLMs (n/a represents not available.). ToolACE comprehensively incorporates the broadest range of APIs and domains, supports complex nested parameters (Nested), accommodates both parallel (Parallel) and dependent (Dependent) function calls, and addresses various types of tool-related data (Multi-type).}
\label{table:comparison}
\vspace{-0.4cm}
\begin{tabular}{ccccccc}
\toprule
Model      & \#API & \#Domain & Nested & Parallel & Dependent & Multi-type \\\midrule
Gorilla~\cite{patil2023gorilla}    & 1645  & 3        & \xmark         & \xmark           &  \xmark              & \xmark        \\
ToolAlpaca~\cite{tang2023toolalpaca} & 3938  & 50       & \xmark         & \xmark          & \xmark             & \xmark         \\
ToolLLM~\cite{qin2023toolllm}    & 16464 & 49   & \xmark    & \xmark         &  \cmark                     & \xmark           \\
Functionary~\cite{functionarymeetkai}    & n/a  & n/a    & \xmark    & \cmark         &  \xmark                     & \xmark           \\
xLAM~\cite{liu2024apigen}       & 3673  & 21    & \xmark   & \cmark         &  \xmark                     & \xmark           \\
Granite~\cite{abdelaziz2024granite}       & n/a  & n/a     & \xmark   & \cmark         &  \xmark                     & \cmark           \\
\textbf{ToolACE} & \textbf{26507} & \textbf{390} & \cmark & \cmark & \cmark & \cmark\\\bottomrule
\end{tabular}
\end{threeparttable}
\end{table}



In this paper, we present ToolACE, a systematic tool-learning pipeline that automatically synthesizes \textit{accurate}, \textit{diverse}, and \textit{complex} function calling data, with the awareness of the model's capability. 

\noindent
\textbf{Evolutionary Diversity.} 
Exposing LLMs to a broad range of function-calling scenarios enhances their proficiency and zero-shot tool usage~\cite{zhang2024instruction}. Instead of relying on public APIs, ToolACE introduces a Tool Self-Evolution Synthesis (TSS) method. TSS uses a speciation-adaptation-evolution process to generate tools across multiple domains with diverse data types and constraints. Starting with pretraining data to ensure comprehensive coverage, this iterative process of self-evolution and continual updates expands the diversity of the API pool, enabling more sophisticated data generation.
% \paragraph{Evolutional Diversity.} 
% Exposing LLMs to diverse function-calling scenarios facilitates a more well-rounded cognitive skill set of tool usage~\cite{zhang2024instruction}. Rather than depending on public APIs, ToolACE introduces a tool self-evolution synthesis (TSS) method, which generates tools covering a wide range of domains, incorporating diverse data types and constraints.
% Based on the pretraining data, an iterative process of self-evolution and updates progressively increases
% the diversity of our API pool for data generation.

\noindent
\textbf{Self-Guided Complexity.} Instruction-following data should possess sufficient complexity to foster function-calling skills. LLMs learn more effectively when the complexity of the data slightly exceeds their current capability~\cite{du2023mods}. To address this, we propose a self-guided dialog generation process (SDG), where the LLM serves as an evaluator to regulate complexity. 
Four types of function-calling data are generated with multi-agent interactions, following a self-guided complication strategy.

\noindent
\textbf{Refined Accuracy.} Data accuracy is fundamental to the effectiveness of tool-augmented LLMs. 
% To produce high-quality data, 
ToolACE employs a dual-layer verification (DLV) system, integrating both rule-based and model-based checkers, to guarantee the executability and consistency of the synthesized data.

Equipped with data accuracy, complexity, and diversity, ToolACE aims to enhance the function-calling capability of LLMs with strong generalization. Our contributions are outlined as follows:
\begin{itemize}[left=0pt]
    \item We propose a novel automated data pipeline for function calls, ToolACE, which comprises a tool self-evolution synthesis module, a self-guided dialog generation module, and a dual-layer verification module. To our knowledge, this is the first work to highlight the benefits of synthesizing diverse APIs to improve the generalization of function calls.
    \item We develop a self-guided complication strategy to generate various types of function-calling dialogs with appropriate complexity. The given LLM is utilized as the complexity evaluator to guide the complexity level of the generated data. The quality of the generated data is ensured through a dual-layer verification process, which combines both rule checkers and model checkers.
    \item We conduct experiments on two widely adopted benchmarks: BFCL~\cite{berkeley-function-calling-leaderboard} and APIBank~\cite{li2023api}. With only 8B parameters, ToolACE significantly outperforms existing open-source LLMs and is competitive with the latest GPT-4 models.
\end{itemize}






% % In ToolACE, we introduce a concept of proximal complexity (PC),
% Data that is either too simple or too complex proves unproductive for LLMs. 
% This aligns with the Zone of Proximal Development (ZPD) theory in educational psychology, which states that learning is most effective when the tasks are just beyond the learner's current level of independent capability but achievable with suitable support~\cite{shabani2010vygotsky}.

% \smallskip
% \noindent




% In this system, a labeling agent evaluates the data's validity based on predefined rules and criteria, which is further verified through spot checks by human experts.




% For instance, when a user requests to book a flight to Denver, instead of providing simple step-by-step instructions, function-calling LLMs can directly search for flight information and assist the user in booking the flight.

% Tool-augmented large language models
% suggest practical recipes for enhancement
% programatically generated data

% \input{sec/further_pretrain}
\input{sec/data_pipeline}
\input{sec/data_verification}
% \input{sec/data_analysis}
\input{sec/experiment}
\input{sec/related_work}

\section{Conclusion}
This paper presents ToolACE, an automated data generation pipeline developed to enhance the function-calling capabilities of large language models. ToolACE employs a novel self-evolution synthesis process and a self-guided data generation method to curate accurate, complex, and diverse synthetic APIs and dialogs. Our results demonstrate that even smaller models trained with ToolACE can achieve state-of-the-art performance, thereby advancing the field and setting new benchmarks for tool-augmented AI agents.



% \section{Acknowledgements}
% We thank the contribution of Mei Li, Wei Tang, Xi Wang, Meide Zhang, Ke Xiong, and Rencong Shi to this work.

% rivaling advanced models like GPT-4. 
% \newpage

\bibliographystyle{iclr2025_conference}
\bibliography{ref}

\input{sec/appendix}


\end{document}
